% =====================================================================
% Section 1 -- Introduction
% Usage: % =====================================================================
% Section 1 -- Introduction
% Usage: % =====================================================================
% Section 1 -- Introduction
% Usage: % =====================================================================
% Section 1 -- Introduction
% Usage: \input{sec1_intro} inside an acmart document (see main.tex).
% =====================================================================
\providecommand{\TBD}[1][]{\textcolor{red}{[TBD#1]}}

\section{Introduction}
\label{sec:intro}

Every smartphone that rides in a car carries an accelerometer, a gyroscope and a GNSS receiver.
Two decades of mobile sensing have used them to monitor roads: early systems detected potholes and braking from phones and dedicated in-vehicle units~\cite{mohan2008nericell,eriksson2008pothole}, and a large body of later work classifies surface types or regresses a roughness index from vehicle vibration~\cite{menegazzo2021pvs}.
The output of these systems is a label or a number per window or per road segment.
The people who drive the roads ask something different.
A commuter wants to know which part of the route was rough and whether to drive it differently.
A passenger prone to car sickness wants to know which stretches and which driving made the ride unpleasant.
A fleet manager wants to know where the vans take shocks, and whether the map is wrong about a surface.
These are questions about a specific trip, in language, and their answers combine facts that a single number cannot hold.

\paragraph{Why a scalar roughness score is the wrong output.}
The International Roughness Index (IRI) is defined on a measured road profile, through a reference quarter-car model driven at a fixed speed~\cite{sayers1986iri}.
A phone in a holder observes something else: the response of an unknown vehicle, through an unknown mount, at whatever speed the driver chose.
Regressing IRI, or any single roughness score, from such data runs into four problems.
\begin{enumerate}
\item \emph{It confounds the road with the drive.} Vibration grows with speed, so a regressor can score well by learning speed. On the PVS dataset~\cite{menegazzo2021pvs}, speed alone separates the three surface types with an AUC of 0.977 \TBD[: re-run audit], and an earlier ``kinetic normalization'' of our own encoded speed almost perfectly (Spearman 0.995 with speed) \TBD[: re-run audit].
\item \emph{It collapses distinct facts.} How the ride \emph{felt} depends on speed; how the \emph{road} is should not. A stretch with one violent speed bump and a stretch of steady cobblestone can share a mean vibration level, yet ISO~2631-1 treats the first as shock-dominated and requires different measures for it~\cite{iso2631}.
\item \emph{It has no place, cause or confidence.} A number does not say where along the street the roughness was, whether a jolt came from the road or from pulling away, or how sure the system is.
\item \emph{It cannot say what is missing.} When the GNSS fix drops out, the gyroscope is absent or the map contradicts the signal, a score is silently wrong; a sentence can say so.
\end{enumerate}

\paragraph{Why not a language model on the raw signal, or a template?}
Large language models (LLMs) can describe and reason about inertial data~\cite{xu2024penetrative,li2025sensorllm,imran2024llasa,yoon2024bymyeyes}, and deployed driver-coaching systems already turn phone telematics into text~\cite{braun2018saferdrive,baig2022drivingbeacon,thompson2026localising}.
Neither line of work is sufficient here.
Sensor--LLM systems interpret activity windows; they do not segment a trip along the road network, and no claim they make is checked against the measurements.
The coaching systems keep content selection in rules and let a language model at most paraphrase~\cite{thompson2026localising}; they do not cover vibration or road surface, and they do not evaluate the factuality of the generated text.
A carefully engineered template narrator is faithful by construction and is a strong baseline~\cite{vandeemter2005templates}, but it can only produce what its designers enumerated.
A trip capsule in our system can take on the order of $10^4$ joint states per road unit before any question, audience or history enters (Section~\ref{sec:related}).

\paragraph{Our position.}
We do not claim that a small language model (SLM) beats a template narrator at summarizing a fixed schema; on that task a template wins on fidelity and cost.
Our thesis is narrower and testable: \emph{a grounded on-device SLM extends a deterministic narrator to outputs that cannot be enumerated at design time---questions users actually ask, audiences no template was written for, evidence states nobody anticipated---and the fidelity it costs can be measured and bounded.}
RoadScribe therefore keeps the template narrator inside the system as a fallback, never lets the language model supply a value, and measures what the model adds over the template.

\paragraph{RoadScribe.}
RoadScribe turns a trip recorded by a phone into a short narrative and answers follow-up questions about it (Figure~\ref{fig:arch}).
Layer~1 estimates the phone's orientation in the vehicle, computes ISO~2631-1-weighted cabin vibration, separates how the ride felt from how the road is, and detects impulsive events with their maneuver context.
Layer~2 matches the trip to OpenStreetMap~\cite{haklay2008osm} and partitions it into \emph{homogeneous road narrative units} (HRNUs) with a change-point detector whose penalty drops at map transitions.
Layer~3 compiles the units into a typed claim ledger and a token-budgeted capsule, and generates text with an on-device SLM under a plan--ground--realize protocol: the model plans over ledger identifiers, deterministic code fills in every value, and a verifier checks each sentence, falling back to the template narrator when it fails.

\paragraph{Contributions.}
The paper makes three novel technical contributions.
\begin{itemize}
\item \textbf{C1 --- Sensing: map-anchored, uncertainty-annotated ride facts from a phone.}
  A Layer~1 pipeline estimates the mount orientation per handling episode with a gated, rotation-equivariant gravity estimate and a closed-form Procrustes heading fit against GNSS; applies the ISO~2631-1 weightings exactly in the frequency domain; reports felt and speed-decoupled road vibration separately with the standard's shock-evaluation rule; and labels impulsive events with a maneuver gate grounded in measured evidence.
  On PVS the heading estimate recovers a rigid mount to within about $2^\circ$ from trip to trip and exposes a $14^\circ$ mounting offset in one vehicle (Section~\ref{sec:l1eval}).
\item \textbf{C2 --- Segmentation: Map-Informed PELT with a corrected pruning rule.}
  We formulate trip segmentation as penalized change-point detection whose penalty is lowered near OpenStreetMap transitions.
  The classical PELT pruning test~\cite{killick2012pelt} assumes a constant penalty; we prove a three-line lemma giving the correct test, $G(t)+C(\mathbf y_{t+1:s})\ge G(s)$, for location-dependent penalties, and extend it to a minimum segment length.
  In 300 random series with map-dependent penalties, the textbook rule returned a sub-optimal partition in 136, while the corrected rule matched exhaustive optimal partitioning in all 300 (Section~\ref{sec:layer2}).
\item \textbf{C3 --- Generation: Plan--Ground--Realize with a verifier and template fallback.}
  The SLM chooses and orders content by emitting a plan constrained by a grammar generated from the capsule's ledger, so it can reference only units and facts that exist; grounding is deterministic; a claim verifier checks every sentence against the ledger, including events, hedges and fields marked missing; and a sentence that fails twice is replaced by the template realization of the same plan item.
  The language model therefore controls selection, ordering, aggregation and phrasing, but not values (Section~\ref{sec:layer3}).
\end{itemize}
In support of these, we contribute a benchmark whose metrics separate language errors from sensing errors (claim precision against the ledger and against ground truth, kinetic--semantic fidelity, anomaly grounding), and the design of a preregistered within-subject study with 28 drivers plus a two-week field deployment to measure usefulness, trust, workload, situation awareness and error detection against strong template baselines (Section~\ref{sec:eval}).
The implementation, with tests that reproduce every number we report as verified, is released as open source.

\paragraph{Results in brief.}
\TBD[: one paragraph with the headline results of E1--E5 and H1--H4 once the benchmark and the study have run. Until then, no result may be stated here.]

The rest of the paper is organized as follows.
Section~\ref{sec:related} positions RoadScribe against sensor--language models, trajectory--language models and data-to-text driver feedback.
Section~\ref{sec:system} formalizes the three layers.
Section~\ref{sec:eval} describes the data, metrics, benchmark and user study.
Section~\ref{sec:discussion} discusses limitations and ethics.
 inside an acmart document (see main.tex).
% =====================================================================
\providecommand{\TBD}[1][]{\textcolor{red}{[TBD#1]}}

\section{Introduction}
\label{sec:intro}

Every smartphone that rides in a car carries an accelerometer, a gyroscope and a GNSS receiver.
Two decades of mobile sensing have used them to monitor roads: early systems detected potholes and braking from phones and dedicated in-vehicle units~\cite{mohan2008nericell,eriksson2008pothole}, and a large body of later work classifies surface types or regresses a roughness index from vehicle vibration~\cite{menegazzo2021pvs}.
The output of these systems is a label or a number per window or per road segment.
The people who drive the roads ask something different.
A commuter wants to know which part of the route was rough and whether to drive it differently.
A passenger prone to car sickness wants to know which stretches and which driving made the ride unpleasant.
A fleet manager wants to know where the vans take shocks, and whether the map is wrong about a surface.
These are questions about a specific trip, in language, and their answers combine facts that a single number cannot hold.

\paragraph{Why a scalar roughness score is the wrong output.}
The International Roughness Index (IRI) is defined on a measured road profile, through a reference quarter-car model driven at a fixed speed~\cite{sayers1986iri}.
A phone in a holder observes something else: the response of an unknown vehicle, through an unknown mount, at whatever speed the driver chose.
Regressing IRI, or any single roughness score, from such data runs into four problems.
\begin{enumerate}
\item \emph{It confounds the road with the drive.} Vibration grows with speed, so a regressor can score well by learning speed. On the PVS dataset~\cite{menegazzo2021pvs}, speed alone separates the three surface types with an AUC of 0.977 \TBD[: re-run audit], and an earlier ``kinetic normalization'' of our own encoded speed almost perfectly (Spearman 0.995 with speed) \TBD[: re-run audit].
\item \emph{It collapses distinct facts.} How the ride \emph{felt} depends on speed; how the \emph{road} is should not. A stretch with one violent speed bump and a stretch of steady cobblestone can share a mean vibration level, yet ISO~2631-1 treats the first as shock-dominated and requires different measures for it~\cite{iso2631}.
\item \emph{It has no place, cause or confidence.} A number does not say where along the street the roughness was, whether a jolt came from the road or from pulling away, or how sure the system is.
\item \emph{It cannot say what is missing.} When the GNSS fix drops out, the gyroscope is absent or the map contradicts the signal, a score is silently wrong; a sentence can say so.
\end{enumerate}

\paragraph{Why not a language model on the raw signal, or a template?}
Large language models (LLMs) can describe and reason about inertial data~\cite{xu2024penetrative,li2025sensorllm,imran2024llasa,yoon2024bymyeyes}, and deployed driver-coaching systems already turn phone telematics into text~\cite{braun2018saferdrive,baig2022drivingbeacon,thompson2026localising}.
Neither line of work is sufficient here.
Sensor--LLM systems interpret activity windows; they do not segment a trip along the road network, and no claim they make is checked against the measurements.
The coaching systems keep content selection in rules and let a language model at most paraphrase~\cite{thompson2026localising}; they do not cover vibration or road surface, and they do not evaluate the factuality of the generated text.
A carefully engineered template narrator is faithful by construction and is a strong baseline~\cite{vandeemter2005templates}, but it can only produce what its designers enumerated.
A trip capsule in our system can take on the order of $10^4$ joint states per road unit before any question, audience or history enters (Section~\ref{sec:related}).

\paragraph{Our position.}
We do not claim that a small language model (SLM) beats a template narrator at summarizing a fixed schema; on that task a template wins on fidelity and cost.
Our thesis is narrower and testable: \emph{a grounded on-device SLM extends a deterministic narrator to outputs that cannot be enumerated at design time---questions users actually ask, audiences no template was written for, evidence states nobody anticipated---and the fidelity it costs can be measured and bounded.}
RoadScribe therefore keeps the template narrator inside the system as a fallback, never lets the language model supply a value, and measures what the model adds over the template.

\paragraph{RoadScribe.}
RoadScribe turns a trip recorded by a phone into a short narrative and answers follow-up questions about it (Figure~\ref{fig:arch}).
Layer~1 estimates the phone's orientation in the vehicle, computes ISO~2631-1-weighted cabin vibration, separates how the ride felt from how the road is, and detects impulsive events with their maneuver context.
Layer~2 matches the trip to OpenStreetMap~\cite{haklay2008osm} and partitions it into \emph{homogeneous road narrative units} (HRNUs) with a change-point detector whose penalty drops at map transitions.
Layer~3 compiles the units into a typed claim ledger and a token-budgeted capsule, and generates text with an on-device SLM under a plan--ground--realize protocol: the model plans over ledger identifiers, deterministic code fills in every value, and a verifier checks each sentence, falling back to the template narrator when it fails.

\paragraph{Contributions.}
The paper makes three novel technical contributions.
\begin{itemize}
\item \textbf{C1 --- Sensing: map-anchored, uncertainty-annotated ride facts from a phone.}
  A Layer~1 pipeline estimates the mount orientation per handling episode with a gated, rotation-equivariant gravity estimate and a closed-form Procrustes heading fit against GNSS; applies the ISO~2631-1 weightings exactly in the frequency domain; reports felt and speed-decoupled road vibration separately with the standard's shock-evaluation rule; and labels impulsive events with a maneuver gate grounded in measured evidence.
  On PVS the heading estimate recovers a rigid mount to within about $2^\circ$ from trip to trip and exposes a $14^\circ$ mounting offset in one vehicle (Section~\ref{sec:l1eval}).
\item \textbf{C2 --- Segmentation: Map-Informed PELT with a corrected pruning rule.}
  We formulate trip segmentation as penalized change-point detection whose penalty is lowered near OpenStreetMap transitions.
  The classical PELT pruning test~\cite{killick2012pelt} assumes a constant penalty; we prove a three-line lemma giving the correct test, $G(t)+C(\mathbf y_{t+1:s})\ge G(s)$, for location-dependent penalties, and extend it to a minimum segment length.
  In 300 random series with map-dependent penalties, the textbook rule returned a sub-optimal partition in 136, while the corrected rule matched exhaustive optimal partitioning in all 300 (Section~\ref{sec:layer2}).
\item \textbf{C3 --- Generation: Plan--Ground--Realize with a verifier and template fallback.}
  The SLM chooses and orders content by emitting a plan constrained by a grammar generated from the capsule's ledger, so it can reference only units and facts that exist; grounding is deterministic; a claim verifier checks every sentence against the ledger, including events, hedges and fields marked missing; and a sentence that fails twice is replaced by the template realization of the same plan item.
  The language model therefore controls selection, ordering, aggregation and phrasing, but not values (Section~\ref{sec:layer3}).
\end{itemize}
In support of these, we contribute a benchmark whose metrics separate language errors from sensing errors (claim precision against the ledger and against ground truth, kinetic--semantic fidelity, anomaly grounding), and the design of a preregistered within-subject study with 28 drivers plus a two-week field deployment to measure usefulness, trust, workload, situation awareness and error detection against strong template baselines (Section~\ref{sec:eval}).
The implementation, with tests that reproduce every number we report as verified, is released as open source.

\paragraph{Results in brief.}
\TBD[: one paragraph with the headline results of E1--E5 and H1--H4 once the benchmark and the study have run. Until then, no result may be stated here.]

The rest of the paper is organized as follows.
Section~\ref{sec:related} positions RoadScribe against sensor--language models, trajectory--language models and data-to-text driver feedback.
Section~\ref{sec:system} formalizes the three layers.
Section~\ref{sec:eval} describes the data, metrics, benchmark and user study.
Section~\ref{sec:discussion} discusses limitations and ethics.
 inside an acmart document (see main.tex).
% =====================================================================
\providecommand{\TBD}[1][]{\textcolor{red}{[TBD#1]}}

\section{Introduction}
\label{sec:intro}

Every smartphone that rides in a car carries an accelerometer, a gyroscope and a GNSS receiver.
Two decades of mobile sensing have used them to monitor roads: early systems detected potholes and braking from phones and dedicated in-vehicle units~\cite{mohan2008nericell,eriksson2008pothole}, and a large body of later work classifies surface types or regresses a roughness index from vehicle vibration~\cite{menegazzo2021pvs}.
The output of these systems is a label or a number per window or per road segment.
The people who drive the roads ask something different.
A commuter wants to know which part of the route was rough and whether to drive it differently.
A passenger prone to car sickness wants to know which stretches and which driving made the ride unpleasant.
A fleet manager wants to know where the vans take shocks, and whether the map is wrong about a surface.
These are questions about a specific trip, in language, and their answers combine facts that a single number cannot hold.

\paragraph{Why a scalar roughness score is the wrong output.}
The International Roughness Index (IRI) is defined on a measured road profile, through a reference quarter-car model driven at a fixed speed~\cite{sayers1986iri}.
A phone in a holder observes something else: the response of an unknown vehicle, through an unknown mount, at whatever speed the driver chose.
Regressing IRI, or any single roughness score, from such data runs into four problems.
\begin{enumerate}
\item \emph{It confounds the road with the drive.} Vibration grows with speed, so a regressor can score well by learning speed. On the PVS dataset~\cite{menegazzo2021pvs}, speed alone separates the three surface types with an AUC of 0.977 \TBD[: re-run audit], and an earlier ``kinetic normalization'' of our own encoded speed almost perfectly (Spearman 0.995 with speed) \TBD[: re-run audit].
\item \emph{It collapses distinct facts.} How the ride \emph{felt} depends on speed; how the \emph{road} is should not. A stretch with one violent speed bump and a stretch of steady cobblestone can share a mean vibration level, yet ISO~2631-1 treats the first as shock-dominated and requires different measures for it~\cite{iso2631}.
\item \emph{It has no place, cause or confidence.} A number does not say where along the street the roughness was, whether a jolt came from the road or from pulling away, or how sure the system is.
\item \emph{It cannot say what is missing.} When the GNSS fix drops out, the gyroscope is absent or the map contradicts the signal, a score is silently wrong; a sentence can say so.
\end{enumerate}

\paragraph{Why not a language model on the raw signal, or a template?}
Large language models (LLMs) can describe and reason about inertial data~\cite{xu2024penetrative,li2025sensorllm,imran2024llasa,yoon2024bymyeyes}, and deployed driver-coaching systems already turn phone telematics into text~\cite{braun2018saferdrive,baig2022drivingbeacon,thompson2026localising}.
Neither line of work is sufficient here.
Sensor--LLM systems interpret activity windows; they do not segment a trip along the road network, and no claim they make is checked against the measurements.
The coaching systems keep content selection in rules and let a language model at most paraphrase~\cite{thompson2026localising}; they do not cover vibration or road surface, and they do not evaluate the factuality of the generated text.
A carefully engineered template narrator is faithful by construction and is a strong baseline~\cite{vandeemter2005templates}, but it can only produce what its designers enumerated.
A trip capsule in our system can take on the order of $10^4$ joint states per road unit before any question, audience or history enters (Section~\ref{sec:related}).

\paragraph{Our position.}
We do not claim that a small language model (SLM) beats a template narrator at summarizing a fixed schema; on that task a template wins on fidelity and cost.
Our thesis is narrower and testable: \emph{a grounded on-device SLM extends a deterministic narrator to outputs that cannot be enumerated at design time---questions users actually ask, audiences no template was written for, evidence states nobody anticipated---and the fidelity it costs can be measured and bounded.}
RoadScribe therefore keeps the template narrator inside the system as a fallback, never lets the language model supply a value, and measures what the model adds over the template.

\paragraph{RoadScribe.}
RoadScribe turns a trip recorded by a phone into a short narrative and answers follow-up questions about it (Figure~\ref{fig:arch}).
Layer~1 estimates the phone's orientation in the vehicle, computes ISO~2631-1-weighted cabin vibration, separates how the ride felt from how the road is, and detects impulsive events with their maneuver context.
Layer~2 matches the trip to OpenStreetMap~\cite{haklay2008osm} and partitions it into \emph{homogeneous road narrative units} (HRNUs) with a change-point detector whose penalty drops at map transitions.
Layer~3 compiles the units into a typed claim ledger and a token-budgeted capsule, and generates text with an on-device SLM under a plan--ground--realize protocol: the model plans over ledger identifiers, deterministic code fills in every value, and a verifier checks each sentence, falling back to the template narrator when it fails.

\paragraph{Contributions.}
The paper makes three novel technical contributions.
\begin{itemize}
\item \textbf{C1 --- Sensing: map-anchored, uncertainty-annotated ride facts from a phone.}
  A Layer~1 pipeline estimates the mount orientation per handling episode with a gated, rotation-equivariant gravity estimate and a closed-form Procrustes heading fit against GNSS; applies the ISO~2631-1 weightings exactly in the frequency domain; reports felt and speed-decoupled road vibration separately with the standard's shock-evaluation rule; and labels impulsive events with a maneuver gate grounded in measured evidence.
  On PVS the heading estimate recovers a rigid mount to within about $2^\circ$ from trip to trip and exposes a $14^\circ$ mounting offset in one vehicle (Section~\ref{sec:l1eval}).
\item \textbf{C2 --- Segmentation: Map-Informed PELT with a corrected pruning rule.}
  We formulate trip segmentation as penalized change-point detection whose penalty is lowered near OpenStreetMap transitions.
  The classical PELT pruning test~\cite{killick2012pelt} assumes a constant penalty; we prove a three-line lemma giving the correct test, $G(t)+C(\mathbf y_{t+1:s})\ge G(s)$, for location-dependent penalties, and extend it to a minimum segment length.
  In 300 random series with map-dependent penalties, the textbook rule returned a sub-optimal partition in 136, while the corrected rule matched exhaustive optimal partitioning in all 300 (Section~\ref{sec:layer2}).
\item \textbf{C3 --- Generation: Plan--Ground--Realize with a verifier and template fallback.}
  The SLM chooses and orders content by emitting a plan constrained by a grammar generated from the capsule's ledger, so it can reference only units and facts that exist; grounding is deterministic; a claim verifier checks every sentence against the ledger, including events, hedges and fields marked missing; and a sentence that fails twice is replaced by the template realization of the same plan item.
  The language model therefore controls selection, ordering, aggregation and phrasing, but not values (Section~\ref{sec:layer3}).
\end{itemize}
In support of these, we contribute a benchmark whose metrics separate language errors from sensing errors (claim precision against the ledger and against ground truth, kinetic--semantic fidelity, anomaly grounding), and the design of a preregistered within-subject study with 28 drivers plus a two-week field deployment to measure usefulness, trust, workload, situation awareness and error detection against strong template baselines (Section~\ref{sec:eval}).
The implementation, with tests that reproduce every number we report as verified, is released as open source.

\paragraph{Results in brief.}
\TBD[: one paragraph with the headline results of E1--E5 and H1--H4 once the benchmark and the study have run. Until then, no result may be stated here.]

The rest of the paper is organized as follows.
Section~\ref{sec:related} positions RoadScribe against sensor--language models, trajectory--language models and data-to-text driver feedback.
Section~\ref{sec:system} formalizes the three layers.
Section~\ref{sec:eval} describes the data, metrics, benchmark and user study.
Section~\ref{sec:discussion} discusses limitations and ethics.
 inside an acmart document (see main.tex).
% =====================================================================
\providecommand{\TBD}[1][]{\textcolor{red}{[TBD#1]}}

\section{Introduction}
\label{sec:intro}

Every smartphone that rides in a car carries an accelerometer, a gyroscope and a GNSS receiver.
Two decades of mobile sensing have used them to monitor roads: early systems detected potholes and braking from phones and dedicated in-vehicle units~\cite{mohan2008nericell,eriksson2008pothole}, and a large body of later work classifies surface types or regresses a roughness index from vehicle vibration~\cite{menegazzo2021pvs}.
The output of these systems is a label or a number per window or per road segment.
The people who drive the roads ask something different.
A commuter wants to know which part of the route was rough and whether to drive it differently.
A passenger prone to car sickness wants to know which stretches and which driving made the ride unpleasant.
A fleet manager wants to know where the vans take shocks, and whether the map is wrong about a surface.
These are questions about a specific trip, in language, and their answers combine facts that a single number cannot hold.

\paragraph{Why a scalar roughness score is the wrong output.}
The International Roughness Index (IRI) is defined on a measured road profile, through a reference quarter-car model driven at a fixed speed~\cite{sayers1986iri}.
A phone in a holder observes something else: the response of an unknown vehicle, through an unknown mount, at whatever speed the driver chose.
Regressing IRI, or any single roughness score, from such data runs into four problems.
\begin{enumerate}
\item \emph{It confounds the road with the drive.} Vibration grows with speed, so a regressor can score well by learning speed. On the PVS dataset~\cite{menegazzo2021pvs}, speed alone separates the three surface types with an AUC of 0.977 \TBD[: re-run audit], and an earlier ``kinetic normalization'' of our own encoded speed almost perfectly (Spearman 0.995 with speed) \TBD[: re-run audit].
\item \emph{It collapses distinct facts.} How the ride \emph{felt} depends on speed; how the \emph{road} is should not. A stretch with one violent speed bump and a stretch of steady cobblestone can share a mean vibration level, yet ISO~2631-1 treats the first as shock-dominated and requires different measures for it~\cite{iso2631}.
\item \emph{It has no place, cause or confidence.} A number does not say where along the street the roughness was, whether a jolt came from the road or from pulling away, or how sure the system is.
\item \emph{It cannot say what is missing.} When the GNSS fix drops out, the gyroscope is absent or the map contradicts the signal, a score is silently wrong; a sentence can say so.
\end{enumerate}

\paragraph{Why not a language model on the raw signal, or a template?}
Large language models (LLMs) can describe and reason about inertial data~\cite{xu2024penetrative,li2025sensorllm,imran2024llasa,yoon2024bymyeyes}, and deployed driver-coaching systems already turn phone telematics into text~\cite{braun2018saferdrive,baig2022drivingbeacon,thompson2026localising}.
Neither line of work is sufficient here.
Sensor--LLM systems interpret activity windows; they do not segment a trip along the road network, and no claim they make is checked against the measurements.
The coaching systems keep content selection in rules and let a language model at most paraphrase~\cite{thompson2026localising}; they do not cover vibration or road surface, and they do not evaluate the factuality of the generated text.
A carefully engineered template narrator is faithful by construction and is a strong baseline~\cite{vandeemter2005templates}, but it can only produce what its designers enumerated.
A trip capsule in our system can take on the order of $10^4$ joint states per road unit before any question, audience or history enters (Section~\ref{sec:related}).

\paragraph{Our position.}
We do not claim that a small language model (SLM) beats a template narrator at summarizing a fixed schema; on that task a template wins on fidelity and cost.
Our thesis is narrower and testable: \emph{a grounded on-device SLM extends a deterministic narrator to outputs that cannot be enumerated at design time---questions users actually ask, audiences no template was written for, evidence states nobody anticipated---and the fidelity it costs can be measured and bounded.}
RoadScribe therefore keeps the template narrator inside the system as a fallback, never lets the language model supply a value, and measures what the model adds over the template.

\paragraph{RoadScribe.}
RoadScribe turns a trip recorded by a phone into a short narrative and answers follow-up questions about it (Figure~\ref{fig:arch}).
Layer~1 estimates the phone's orientation in the vehicle, computes ISO~2631-1-weighted cabin vibration, separates how the ride felt from how the road is, and detects impulsive events with their maneuver context.
Layer~2 matches the trip to OpenStreetMap~\cite{haklay2008osm} and partitions it into \emph{homogeneous road narrative units} (HRNUs) with a change-point detector whose penalty drops at map transitions.
Layer~3 compiles the units into a typed claim ledger and a token-budgeted capsule, and generates text with an on-device SLM under a plan--ground--realize protocol: the model plans over ledger identifiers, deterministic code fills in every value, and a verifier checks each sentence, falling back to the template narrator when it fails.

\paragraph{Contributions.}
The paper makes three novel technical contributions.
\begin{itemize}
\item \textbf{C1 --- Sensing: map-anchored, uncertainty-annotated ride facts from a phone.}
  A Layer~1 pipeline estimates the mount orientation per handling episode with a gated, rotation-equivariant gravity estimate and a closed-form Procrustes heading fit against GNSS; applies the ISO~2631-1 weightings exactly in the frequency domain; reports felt and speed-decoupled road vibration separately with the standard's shock-evaluation rule; and labels impulsive events with a maneuver gate grounded in measured evidence.
  On PVS the heading estimate recovers a rigid mount to within about $2^\circ$ from trip to trip and exposes a $14^\circ$ mounting offset in one vehicle (Section~\ref{sec:l1eval}).
\item \textbf{C2 --- Segmentation: Map-Informed PELT with a corrected pruning rule.}
  We formulate trip segmentation as penalized change-point detection whose penalty is lowered near OpenStreetMap transitions.
  The classical PELT pruning test~\cite{killick2012pelt} assumes a constant penalty; we prove a three-line lemma giving the correct test, $G(t)+C(\mathbf y_{t+1:s})\ge G(s)$, for location-dependent penalties, and extend it to a minimum segment length.
  In 300 random series with map-dependent penalties, the textbook rule returned a sub-optimal partition in 136, while the corrected rule matched exhaustive optimal partitioning in all 300 (Section~\ref{sec:layer2}).
\item \textbf{C3 --- Generation: Plan--Ground--Realize with a verifier and template fallback.}
  The SLM chooses and orders content by emitting a plan constrained by a grammar generated from the capsule's ledger, so it can reference only units and facts that exist; grounding is deterministic; a claim verifier checks every sentence against the ledger, including events, hedges and fields marked missing; and a sentence that fails twice is replaced by the template realization of the same plan item.
  The language model therefore controls selection, ordering, aggregation and phrasing, but not values (Section~\ref{sec:layer3}).
\end{itemize}
In support of these, we contribute a benchmark whose metrics separate language errors from sensing errors (claim precision against the ledger and against ground truth, kinetic--semantic fidelity, anomaly grounding), and the design of a preregistered within-subject study with 28 drivers plus a two-week field deployment to measure usefulness, trust, workload, situation awareness and error detection against strong template baselines (Section~\ref{sec:eval}).
The implementation, with tests that reproduce every number we report as verified, is released as open source.

\paragraph{Results in brief.}
\TBD[: one paragraph with the headline results of E1--E5 and H1--H4 once the benchmark and the study have run. Until then, no result may be stated here.]

The rest of the paper is organized as follows.
Section~\ref{sec:related} positions RoadScribe against sensor--language models, trajectory--language models and data-to-text driver feedback.
Section~\ref{sec:system} formalizes the three layers.
Section~\ref{sec:eval} describes the data, metrics, benchmark and user study.
Section~\ref{sec:discussion} discusses limitations and ethics.
